\bookchapter{Capítulo 58: Manual de consulta rápida e recursos}

\chapterintro{Este capítulo é um manual de consulta rápida que pode ser usado de forma independente. Quando você não lembrar de alguma sintaxe ou precisar decidir rápido, venha direto para cá. Trate-o como um 'cartão de bolso': consultas de código frequentes, perguntas comuns e recursos de estudo — tudo ali, sem deixar nada de fora.}

\booksection{Cartão de consulta rápida de código}

\begin{center}
\begin{tabularx}{\linewidth}{|X|X|}
\hline
\textbf{O que fazer} & \textbf{Código-chave} \tabularnewline
\hline
Guardar um estado mutável & var x by remember \{ mutableStateOf(valorInicial) \} \tabularnewline
Não perder o estado ao girar a tela & rememberSaveable \{ mutableStateOf(valorInicial) \} \tabularnewline
Layout vertical/horizontal & \texttt{Column \{ \}}/\texttt{Row \{ \}} \tabularnewline
Lista de carregamento preguiçoso & LazyColumn \{ items(lista) \{ ... \} \} \tabularnewline
Navegar entre telas & navController.navigate("rota") \tabularnewline
Operações de espera, dentro de corrotinas & \texttt{viewModelScope.\allowbreak launch \{ .\allowbreak .\allowbreak .\allowbreak  \}} \tabularnewline
Observar fluxos de dados & \texttt{val list by flow.\allowbreak collectAsState()} \tabularnewline
Guardar configurações pequenas & context.dataStore.edit \{ it[KEY] = valor \} \tabularnewline
Guardar dados estruturados & Trio do Room: @Entity / @Dao / @Database \tabularnewline
Buscar dados na rede & Interface Retrofit + funções suspend \tabularnewline
Enviar notificação & NotificationChannel + NotificationCompat.Builder \tabularnewline
Gravar logs & \texttt{Log.\allowbreak d("TAG", "message")} \tabularnewline
\hline
\end{tabularx}
\end{center}

\begin{center}\small\color{gray}Cartão de consulta rápida do desenvolvimento Android (edição de agosto de 2026)\end{center}

\booksection{Perguntas frequentes (FAQ)}

P1: Não sei nada de programação. Quanto tempo leva para aprender desenvolvimento Android?

Siga o caminho deste livro: semanas 1-2, instalar o ambiente e aprender o básico de Kotlin; semanas 3-4, aprender Compose e criar interfaces; semanas 5-6, concluir o projeto do controle de gastos; semanas 7-8, terminar o projeto de clima. Dois meses, do zero a conseguir fazer um App pequeno por conta própria. A condição: dedicar 1 a 2 horas por dia e digitar código de verdade, em vez de só ler.

P2: Ainda preciso aprender Java?

Não precisa. Projetos novos, tutoriais novos e vagas novas praticamente todos usam Kotlin. Você só vai encontrar muito Java ao 'manter um projeto com dez anos de idade'. Depois de aprender Kotlin, se o trabalho pedir, completar com Java é rápido — por baixo, os dois vivem no mesmo ecossistema de máquina virtual.

P3: Emulador ou aparelho real?

Para desenvolver e depurar, o emulador é prático (snapshots, vários modelos), mas a aceitação precisa passar por um aparelho real — toque, desempenho, rede de verdade, sensores: o emulador não substitui. Se as condições permitirem, comprar um aparelho intermediário como máquina de testes é o investimento mais vantajoso.

P4: É fácil encontrar trabalho com desenvolvimento Android?

O cenário de 2026: as vagas iniciantes de 'só sabe fazer interface' têm competição acirrada, mas os desenvolvedores que 'entendem Kotlin + dominam arquitetura + têm obra publicada' continuam sendo disputados. Para quem está aprendendo, o melhor currículo é aquele App feito por você que está na loja.

P5: Fiquei travado num erro. E agora?

Os três passos de ouro: (1) pesquise a primeira linha do erro tal como está no StackOverflow/busca; (2) consulte a página correspondente da documentação oficial; (3) se ainda não resolver, embale 'comportamento esperado + comportamento real + erro completo + código relacionado' e envie à IA ou à comunidade. A esmagadora maioria dos problemas se resolve no primeiro passo.

\booksection{Recursos de estudo recomendados}

\begin{center}
\begin{tabularx}{\linewidth}{|X|X|X|}
\hline
\textbf{Tipo} & \textbf{Recurso} & \textbf{Por que recomendar} \tabularnewline
\hline
Oficial & developer.android.com & Documentação oficial do Android; os tutoriais de Compose/Kotlin estão todos lá; autoridade número um \tabularnewline
Oficial & Site oficial do Kotlin (kotlinlang.org) & Documentação oficial da linguagem + playground online \tabularnewline
Oficial & Google Codelabs & Aulas de laboratório guiadas passo a passo; nota máxima em simpatia para iniciantes \tabularnewline
Comunidade & Stack Overflow (tag android) & Primeiro porto de socorro para erros \tabularnewline
Comunidade & Juejin / seção de Android do CSDN & Muitos artigos práticos em chinês; cuidado para separar os novos dos antigos (veja a data primeiro) \tabularnewline
Vídeo & Canal oficial do Google Developers & As apresentações de novidades da conferência anual de desenvolvedores \tabularnewline
Do livro & Parte 7, 'Prática de Programação com IA' (capítulos 37-47) & O método completo para desenvolver Android com ajuda da IA \tabularnewline
\hline
\end{tabularx}
\end{center}

\begin{center}\small\color{gray}Lista de recursos para aprender desenvolvimento Android (agosto de 2026)\end{center}

\booksection{Comandos Gradle de uso frequente}

No dia a dia, os botões bastam, mas a linha de comando às vezes é mais eficiente (especialmente junto com ferramentas de IA ou pipelines de CI):

\begin{center}
\begin{tabularx}{\linewidth}{|X|X|}
\hline
\textbf{Comando} & \textbf{O que faz} \tabularnewline
\hline
\texttt{.\allowbreak /\allowbreak gradlew assembleDebug} & Gerar um APK de depuração (no Windows, tire o ./) \tabularnewline
\texttt{.\allowbreak /\allowbreak gradlew assembleRelease} & Gerar a versão de release (requer assinatura configurada) \tabularnewline
\texttt{.\allowbreak /\allowbreak gradlew test} & Rodar todos os testes de unidade \tabularnewline
\texttt{.\allowbreak /\allowbreak gradlew lint} & Análise estática: encontra bugs em potencial e código fora do padrão \tabularnewline
\texttt{.\allowbreak /\allowbreak gradlew clean} & Limpar os artefatos de build (a versão de compilação do 'reiniciar resolve') \tabularnewline
\texttt{.\allowbreak /\allowbreak gradlew :app: dependencies} & Ver a árvore de dependências (arma poderosa contra conflitos de versão) \tabularnewline
\hline
\end{tabularx}
\end{center}

\begin{center}\small\color{gray}Comandos Gradle de uso frequente\end{center}

\begin{tcbinfo}{Dica: se o build der problema, comece pelo clean}
Aqueles mistérios de 'ontem compilava, hoje não compila mais', seis ou sete vezes em cada dez, são cache de build bagunçado. Um ./gradlew clean seguido de um novo build cura a maior parte dos problemas 'paranormais'. Se não resolver, reinicie o Android Studio — e use esse truque junto com File Invalidate Caches.

\end{tcbinfo}

\booksection{Palavras finais: do primeiro App ao N-ésimo App}

\booksection{Palavras finais: do primeiro App ao N-ésimo App}

Olhando para trás: você saiu do 'não faço ideia do que é uma Activity' para instalar o ambiente, aprender Kotlin, desenhar interfaces com o Compose, conectar o App a banco de dados e rede, e depois testar, empacotar e publicar — nesse caminho, você já percorreu o trecho mais difícil da jornada do zero ao um.

O mundo da tecnologia muda rápido: frameworks se atualizam, ferramentas se substituem e a IA vai assumindo cada vez mais do trabalho de 'digitar'. Mas uma coisa nunca sai de moda: a sua capacidade de pegar uma ideia vaga, decompô-la em passos claros e implementá-la um a um. Ferramentas são amplificadores; o raciocínio é alicerce. O amplificador pode ser trocado; o alicerce não pode tremer.

Você não é mais só "alguém que sabe usar celular". Você pode transformar as ideias da sua cabeça em apps de verdade rodando em 300 milhões de aparelhos Android. A partir de hoje, crie com o teclado — essa é a poesia de ser desenvolvedor Android.

\booksection{Referências}

\begin{enumerate}
  \item Android Developers. (2026). Android Developer Guides and Jetpack Compose Documentation. developer.android.com.
  \item Kotlin Foundation. (2026). Kotlin Language Documentation. kotlinlang.org.
  \item Android Open Source Project. (2026). Android Platform Version History and API Levels.
  \item Google Play. (2026). Play Console Help: Publishing and App Signing.
  \item AndroidX Release Notes. (2026). Room / DataStore / Navigation / WorkManager releases.
  \item Square. (2026). Retrofit Documentation. github.com/square/retrofit.
\end{enumerate}
